\chapter{Ranking Theory}
\label{sec:p15-rationality-5-3}
\index{ranking theory}\index{degree of disbelief}\index{belief revision}

Gabriele Kern-Isberner, Niels Skovgaard-Olsen and Wolfgang Spohn's
2021 chapter distinguishes believing a proposition from assigning it a
degree of confidence.  Its proposed ranking state represents both outright
belief and the firmness with which it is held, and retains a state on which
further revisions can operate.  The chapter's philosophical case, mathematical
definitions, Tweety examples, AI applications and psychological research
programme are different kinds of claims; a formal representation theorem is
not a demonstration that people use that representation
\parencite[pp.~337--345]{KernIsbernerSkovgaardOlsenSpohn2021Ranking}.
The source uses \emph{doxastic} for belief, reserving the stricter sense of
\emph{epistemic} for knowledge; its AI discussion cautions that those words
are often less sharply separated in knowledge-representation practice.

\section{Why the source wants a theory of outright belief}
\label{sec:p15-w02b-p15-039-motivation}

The chapter's opening argument has five steps, and each rests on a
different kind of premise.
\begin{enumerate}
\item \emph{Two notions.}  Doxastic attitudes come in degrees, under names
 such as degree of belief, uncertainty or plausibility, and the source's
 own handbook surveys several accounts of them.  Alongside these there is
 belief simpliciter, together with the related notions of acceptance and
 judgment.  This is a conceptual distinction, not yet an argument.
\item \emph{Truth versus action.}  The source treats outright belief as the
 more basic notion where truthful representation of the world is at
 issue: a belief can be true or false, whereas a degree of belief cannot,
 and degrees are tied rather to action.  It adds a critical judgment:
 degree-based accounts struggle to accommodate this point, and it is
 questionable to treat degrees as basic while discounting that difficulty.
 This step is a philosophical premise about the bearers of truth; the
 chapter supports it by this contrast, not by a formal result.
\item \emph{Static representation.}  Doxastic logic was the first
 theoretical account of outright belief, but the source calls it static:
 it says what a belief state contains, not how the state changes.
\item \emph{Ordinal dynamics.}  Belief revision theory supplies the
 missing dynamics, but the source says it has problems with iterated
 revision, and that a complete dynamic account requires iteration.  It
 locates the difficulty in the ordinal character of entrenchment orders
 and says that within that framework no clear solution emerged.  It also
 explains why a qualitative belief set cannot do the job by itself: when
 beliefs change, less firmly entrenched beliefs are given up and better
 entrenched ones retained, which calls for an entrenchment order or for
 degrees measuring the strength of entrenchment.
\item \emph{The proposal.}  Ranking theory is offered as representing
 both belief and degrees of belief, and as supplying a complete dynamics
 for both.  This is the claim the remaining sections make precise.
\end{enumerate}
\parencite[pp.~337--338, 341]{KernIsbernerSkovgaardOlsenSpohn2021Ranking}.
Steps 1 and~3 are classifications of existing theories, step~2 is a
contested normative and semantic premise, and step~4 is a historical
assessment of the belief revision literature that the chapter itself later
qualifies: it reports postulates for iterated revision by Darwiche and Pearl
and later authors that were illustrated with ranking functions
\parencite[p.~342]{KernIsbernerSkovgaardOlsenSpohn2021Ranking}.
Step~5 is the formal proposal, and its \emph{completeness} claim means that
every ranking state, after any update of the kind defined in
Definition~\ref{def:p15-rat53-soft-update}, is again a ranking state to which
the same rule applies whenever both new evidence cells have finite rank
(Theorem~\ref{thm:p15-rat53-update}).

The following calculation shows, in a small model, why the source's dynamics
needs more than an order of plausibility.  It is an illustration constructed
for this treatment, not an example printed in the chapter, and it does not show
that every ordinal proposal for iterated revision fails.

\begin{hbexample}{def:p15-w02b-p15-039-same-order}{One plausibility order, two cardinal states, two iterated outcomes}
Let \(W=\{w_1,w_2,w_3\}\) and compare the world ranks
\(r_1=(0,1,2)\) and \(r_2=(0,1,5)\).  Both order the worlds
\(w_1\prec w_2\prec w_3\), so their belief sets coincide at the
thresholds \(0\) and \(1\), and revision by selecting the most plausible
worlds of an input gives the same result for every input.  Now accept
\(E=\{w_2,w_3\}\) with firmness \(1\) by the two-cell rule
of Theorem~\ref{thm:p15-rat53-update}.  In both states \(\kappa(E)=1\) and
\(\kappa(\neg E)=0\), so \(E\)-worlds lose one rank and \(w_1\) gains one:
\[
 r_1'=(1,0,1),\qquad r_2'=(1,0,4).
\]
The two posterior orders differ.  A further acceptance of
\(\neg\{w_2\}=\{w_1,w_3\}\), with any firmness \(n>0\), now selects \(\{w_1,w_3\}\) from
\(r_1'\) but only \(\{w_1\}\) from \(r_2'\).  The difference arises from
the conditional rank \(\kappa(\{w_3\}\mid E)\), which is \(1\) in the first
state and \(4\) in the second.  The order alone does not fix it.
\end{hbexample}

\section{Objects, interpretation and the closure condition}

\begin{hbdefinition}{def:p15-rat53-ranking}{The source's negative, positive and two-sided ranks}
Let \(W\ne\varnothing\) be a space of mutually exclusive, exhaustive
possibilities and identify propositions with subsets of \(W\).
Write \(\neg A=W\setminus A\), \(A\land B=A\cap B\) and
\(A\lor B=A\cup B\), retaining the source's propositional notation.
A negative ranking function is
\(\kappa:\mathcal P(W)\to\mathbb N\cup\{\infty\}\), with
\(0\in\mathbb N\), satisfying
\[
 \kappa(W)=0,\qquad \kappa(\varnothing)=\infty,\qquad
 \kappa(A\lor B)=\min\{\kappa(A),\kappa(B)\}.
\]
Its values measure disbelief, not probabilities: zero means not disbelieved,
not necessarily believed.  Its positive and two-sided counterparts are
\[
 \beta(A)=\kappa(\neg A),\qquad
 \tau(A)=\kappa(\neg A)-\kappa(A).
\]
At least one of the two ranks in the subtraction is zero, so no
\(\infty-\infty\) is used.  Belief, disbelief and suspension are
respectively \(\tau(A)>0\), \(\tau(A)<0\), \(\tau(A)=0\).
The positive version has \(\beta(W)=\infty\),
\(\beta(\varnothing)=0\) and
\(\beta(A\land B)=\min\{\beta(A),\beta(B)\}\).
Conversely \(\kappa(A)=\max\{0,-\tau(A)\}\) and
\(\beta(A)=\max\{0,\tau(A)\}\); these variants carry the same
information, although only the negative ranks obey the displayed
minimum-over-disjunction rule.
For a finite threshold \(z\geq0\), the associated belief set is
\(\mathcal B_z=\{A:\kappa(\neg A)>z\}\); suspension then occupies
the two-sided interval \([-z,z]\).  The source interprets these as
dispositional commitments, not an occurrent thinker consciously entertaining
every consequence \parencite[pp.~337--339]{KernIsbernerSkovgaardOlsenSpohn2021Ranking}.
\end{hbdefinition}

\begin{hbtheorem}{thm:p15-rat53-rank-laws}{proved=annex:p15.rat53.rank-laws}{Belief laws and their finite-space scope}
The event axioms give antitonicity \(A\subseteq B\Rightarrow
\kappa(A)\geq\kappa(B)\), the negation law
\(\min\{\kappa(A),\kappa(\neg A)\}=0\), and finite-conjunction
closure of every \(\mathcal B_z\).
If \(W\) is finite, put \(r(w)=\kappa(\{w\})\).  Then
\[
 \kappa(A)=\min_{w\in A}r(w),\quad \min\varnothing=\infty,
 \qquad
 \mathcal B_z=\{A:M_z\subseteq A\},\quad
 M_z=\{w:r(w)\leq z\}\ne\varnothing.
\]
Hence \(\bigcap\mathcal B_z=M_z\ne\varnothing\), and every event
containing that intersection is believed.  Conversely a normalized
event-valued function whose threshold belief sets are upward closed and
closed under finite intersections at every finite \(z\) satisfies the
binary minimum law.  On an infinite space the same full intersection
conclusion holds if the displayed minimum-over-worlds representation is
assumed with an attained zero minimum; it does not follow from the binary
event axioms alone.
\end{hbtheorem}

\begin{hbexample}{def:p15-rat53-infinite-boundary}{Finite consistency is not a nonempty intersection of all beliefs}
On \(W=\mathbb N\), let \(\kappa(A)=\infty\) for finite \(A\) and
\(\kappa(A)=0\) for infinite \(A\).  The binary source axioms hold.
But \(\mathcal B_0\) is the cofinite filter, whose complete intersection
is empty.  Its finite subfamilies are consistent, and it does not contain
the contradiction.  Thus the source's strong set-theoretic formulation of
consistency and closure needs the finite or world-minimum condition in the
theorem.  A ranking of individual worlds and an event function with only
binary minimum preservation are not automatically the same object on an
infinite domain.
\end{hbexample}

The source treats consistency and closure as rational requirements, while
acknowledging the lottery and preface objections to probabilistic threshold
accounts.  That normative commitment does not follow merely from naming
\(\kappa\).  For a concrete contrast, on ten equiprobable outcomes each
proposition excluding one outcome has probability \(0.9\), but their
conjunction is impossible.  A threshold below \(0.9\) accepts them all;
ranking belief sets under the stated conditions cannot do that.
The source accordingly develops ranking as a theory of outright belief
rather than identifying belief with probability sufficiently close to one
\parencite[pp.~338, 341]{KernIsbernerSkovgaardOlsenSpohn2021Ranking}.

\section{Conditionals and reasons}

The source introduces conditional ranks by a definition, then rewrites the
definition as a law of conjunction,
\(\kappa(A\land B)=\kappa(A)+\kappa(B\mid A)\), and defends that law by an
informal decomposition of how \(A\land B\) can be false.  Either \(A\) is
false, which contributes the disbelief \(\kappa(A)\); or \(A\) is true, in
which case \(B\) must be false for the conjunction to fail, and this adds the
disbelief in \(B\) given \(A\)
\parencite[p.~339]{KernIsbernerSkovgaardOlsenSpohn2021Ranking}.
The argument motivates the definition; it does not derive it from the axioms,
which mention only minima.  Two features of it carry the weight.  First, the
two contributions are \emph{added}, so the definition presupposes that ranks
are cardinal quantities: their differences, not just their order, must be
meaningful.  The source says exactly this, and refers the justification of
cardinality to a measurement result discussed below.  Second, the argument
is sequential, mirroring the probabilistic product rule with addition in
place of multiplication.  The example in
Definition~\ref{def:p15-w02b-p15-039-same-order} shows what the cardinal
reading adds: two states with the same order disagree about a conditional
rank, and therefore about the result of an update.  The chapter's historical
comparison draws the same line, attributing to Shackle, Rescher and Cohen the
minimum-preserving structure but not the conditional ranks that make the
theory cardinal
\parencite[pp.~339, 341]{KernIsbernerSkovgaardOlsenSpohn2021Ranking}.

\begin{hbdefinition}{def:p15-rat53-conditionals}{Conditional ranks and the source's reason relation}
For \(\kappa(A)<\infty\), define
\[
 \kappa(B\mid A)=\kappa(A\land B)-\kappa(A),\qquad
 \beta(B\mid A)=\kappa(\neg B\mid A),\qquad
 \tau(B\mid A)=\kappa(\neg B\mid A)-\kappa(B\mid A).
\]
The condition may be disbelieved: finite positive \(\kappa(A)\) is
permitted.  An infinite-rank condition is outside this subtraction rule.
When both condition cells have finite rank, \(A\) is a reason for \(B\)
relative to this state exactly when
\(\tau(B\mid A)>\tau(B\mid\neg A)\).
The chapter singles out a sufficient reason: \(B\) is believed given
\(A\) and not believed given \(\neg A\).  It compares signed conditional
attitudes, including strengthening acceptance, weakening rejection and
crossing the acceptance boundary; suspension is a separate zero boundary.
If both conditional signs are positive or both are negative, their numerical
ordering is still needed.  A positive sign under \(A\) and a negative sign
under \(\neg A\) gives a sufficient positive reason; reversing those signs
makes \(A\) speak against \(B\), rather than satisfying the strict
positive-reason definition.
Its proposed notion is subject-relative, not a completed account of objective
normative reasons.  Conditional equality
\(\tau(B\mid A)=\tau(B\mid\neg A)\) expresses the local irrelevance
comparison described in the source, not by itself an arbitrary Bayesian-network
factorization theorem \parencite[pp.~339--340]{KernIsbernerSkovgaardOlsenSpohn2021Ranking}.
\end{hbdefinition}

\begin{hbtheorem}{thm:p15-rat53-conditional-laws}{proved=annex:p15.rat53.conditional-laws}{Conditional consistency, finite conditioning and deductive reasons}
For \(\kappa(A)<\infty\), \(B\mapsto\kappa(B\mid A)\) is a
negative ranking function and
\[
 \kappa(A\land B)=\kappa(A)+\kappa(B\mid A),\qquad
 \min\{\kappa(B\mid A),\kappa(\neg B\mid A)\}=0.
\]
If \(A\subseteq B\) and \(\kappa(A),\kappa(\neg B)<\infty\), then
\(A\) is a reason for \(B\) under the displayed definition.  The
conditional consistency law can replace binary minimum preservation in an
axiomatization only with the needed subset-monotonicity condition also present;
normalization and that conditional law alone do not suffice.
\end{hbtheorem}

\subsection{The reason relation and its probabilistic counterpart}
\label{sec:p15-w02b-p15-039-reasons}

The chapter's general explanation of a reason is deliberately prior to any
formalism: \(A\) is a reason for \(B\) when \(A\) speaks in favor of \(B\)
or confirms it, when \(A\) makes \(B\) more plausible or less implausible,
or when \(B\) is more credible or less incredible given \(A\) than given
\(\neg A\).  The source claims that this explanation works for every
conception of conditional degrees of belief.  Read probabilistically, it is
Carnap's incremental confirmation, or positive relevance, which the source
calls basic for confirmation theory; its note~4 refers to Merin's chapter on
relevance and says that the first four sections of that chapter work equally
well with ranking theory.  Definition~\ref{def:p15-rat53-conditionals}
formalizes the last of the listed phrasings, comparing the two-sided rank of
\(B\) given \(A\) with that given \(\neg A\)
\parencite[pp.~339--340 and p.~344, n.~4]{KernIsbernerSkovgaardOlsenSpohn2021Ranking}.
The probabilistic relation is
Definition~\ref{def:p15-r43-polarity} of Chapter~\ref{sec:p15-rationality-4-3}.
The next result tests the note's transfer claim against the polarity laws
of Theorem~\ref{thm:p15-r43-polarity}.  Write the four cell ranks
\[
 a=\kappa(A\land B),\quad b=\kappa(A\land\neg B),\quad
 c=\kappa(\neg A\land B),\quad d=\kappa(\neg A\land\neg B).
\]

\begin{hbtheorem}{thm:p15-w02b-p15-039-reason-laws}{proved=annex:p15-w02b-p15-039-reason-laws}{Cross-sum form of the reason relation and its comparison with positive relevance}
Assume \(\kappa(A),\kappa(\neg A),\kappa(B),\kappa(\neg B)<\infty\).
\begin{enumerate}
\item At most one of the sums \(a+d\) and \(b+c\) is infinite, and \(A\) is
 a reason for \(B\) exactly when \(a+d<b+c\).  Hence \(A\) is a reason for
 \(B\) exactly when \(B\) is a reason for \(A\).
\item \(A\) is a reason for \(\neg B\) exactly when
 \(\tau(B\mid A)<\tau(B\mid\neg A)\), that is, when \(A\) is a reason
 against \(B\); the same holds for \(\neg A\) as a reason for \(B\).
 Neither \(A\) nor \(\neg A\) is a reason for \(B\) exactly when
 \(a+d=b+c\), with both sums finite.
\item For a probability \(P\) with cell masses \(p_1=P(A\land B)\),
 \(p_2=P(A\land\neg B)\), \(p_3=P(\neg A\land B)\), \(p_4=P(\neg A\land\neg B)\),
 the covariance of Definition~\ref{def:p15-r43-polarity} is
 \(\Delta_P(A,B)=p_1p_4-p_2p_3\).  If \(0<P(A)<1\), then
 \(P(B\mid A)>P(B\mid\neg A)\) exactly when \(P(B\mid A)>P(B)\).
\item If all four cell ranks are finite, \(\tau(B)\) lies in the closed
 interval between \(\tau(B\mid A)\) and \(\tau(B\mid\neg A)\).  It need not
 lie strictly inside: for \(a=0,b=1,c=d=5\), \(A\) is a reason for \(B\) and
 \(\tau(B\mid A)=\tau(B)=1\).
\item The reason relation is not transitive.  On \(W=\{1,2,3\}\) with all
 world ranks \(0\), let \(A=\{1\}\), \(C=\{2\}\), \(B=\{1,2\}\).  Then
 \(A\) is a reason for \(B\), \(B\) is a reason for \(C\), and \(A\) is a
 reason against \(C\).
\end{enumerate}
\end{hbtheorem}

Items~1, 2 and~5 confirm the note's transfer claim for the qualitative
polarity laws: symmetry, sign reversal under negation and the failure of
transitivity hold for ranks as for probabilities, and item~3 shows the exact
correspondence, with the product comparison \(p_1p_4\gtrless p_2p_3\) becoming
the sum comparison \(a+d\lessgtr b+c\) (low negative rank meaning high
plausibility).  Item~4 records a material difference.  In probability the
comparison with \(\neg A\) and the comparison with the prior are equivalent
whenever \(0<P(A)<1\); in ranking theory only the former is the source's
definition, and a reason can leave the unconditional two-sided rank unchanged,
because the minimum discards the disbelieved branch.  A ranking version of
Merin's prior-based measures therefore cannot simply be read off from
the source's Definition~3, and the transfer claim of note~4 is a claim about the
qualitative theory; the quantitative measures \(D\), \(Q\) and \(L\) of
Definition~\ref{def:p15-r43-measures} are ratio and difference quantities on
probabilities whose ranking analogues the chapter does not define.
In the Tweety prior of Table~\ref{tab:p15-rat53-tweety}, \(B\) is a reason
for \(F\), since \(0+0<1+0\), and \(P\) is a reason against \(F\), since
\(4+0>1+0\).

The source's Tweety state has three binary properties: bird \(B\), penguin
\(P\), and flight \(F\).  The table below retains all eight numerical
possibilities from its three tables, not just the final favored world.
\begin{table}[htbp]
\centering\small
\begin{tabular}{@{}lrrrr@{}}
\toprule
State and property & \(B\land\neg P\) & \(B\land P\) &
 \(\neg B\land\neg P\) & \(\neg B\land P\)\\
\midrule
Prior \(\kappa\), \(F\) &0&4&0&11\\
Prior \(\kappa\), \(\neg F\)&2&1&0&8\\
After bird evidence \(\kappa'\), \(F\)&0&4&2&13\\
After bird evidence \(\kappa'\), \(\neg F\)&2&1&2&10\\
After penguin evidence \(\kappa''\), \(F\)&1&3&3&12\\
After penguin evidence \(\kappa''\), \(\neg F\)&3&0&3&9\\
\bottomrule
\end{tabular}
\caption{The source's Tweety ranks before and after two specified learning steps.}
\label{tab:p15-rat53-tweety}
\end{table}

Initially the zero-rank worlds satisfy
\(\neg P\land(\neg B\lor F)\).  Tweety is not believed to be a bird,
nor not a bird; flight is also unsettled.  Yet the material implication
\(\neg P\lor F\) is believed because penguinhood is disbelieved.
That does not give belief in flight conditional on penguinhood:
\[
 \kappa(\neg F\mid B)=1,\qquad
 \kappa(\neg F\mid B\land\neg P)=2,\qquad
 \kappa(F\mid P)=3.
\]
The last inequality instead gives conditional belief in nonflight.
This is the chapter's contrast between vacuous material implication and the
Ramsey-test reading of conditional belief
\parencite[Tables~5.3.1--5.3.3, pp.~338--340]{KernIsbernerSkovgaardOlsenSpohn2021Ranking}.

The source generalizes from the example: exactly the material if--then
propositions that are held true non-vacuously correspond to conditional
beliefs.  The vacuous case it has in mind is the penguin conditional, believed
only because the antecedent is disbelieved in the first place
\parencite[pp.~338--339]{KernIsbernerSkovgaardOlsenSpohn2021Ranking}.
The Tweety numbers illustrate the claim but do not establish it.  Reading
``non-vacuously'' as ``with an antecedent that is not disbelieved'', the
claim is the following theorem, and its last item is the limitation.

\begin{hbtheorem}{thm:p15-w02b-p15-039-material}{proved=annex:p15-w02b-p15-039-material}{Material conditionals held non-vacuously are conditional beliefs}
Let \(\kappa(A)<\infty\) and let \(z\geq0\) be a finite belief threshold.
\begin{enumerate}
\item If \(B\) is believed given \(A\) at threshold \(z\), that is
 \(\kappa(\neg B\mid A)>z\), then the material conditional
 \(\neg A\lor B\) is believed at threshold \(z\).
\item If \(\kappa(A)=0\), then \(\neg A\lor B\) is believed at threshold
 \(z\) exactly when \(\kappa(\neg B\mid A)>z\); for \(z=0\) this is
 \(\tau(B\mid A)>0\).
\item If \(\kappa(A)>z\), then \(\neg A\lor B\) is believed at threshold
 \(z\) for every \(B\), whatever the conditional attitude to \(B\) given
 \(A\); in particular \(\neg A\lor B\) and \(\neg A\lor\neg B\) are then
 both believed, while at most one of \(B\) and \(\neg B\) is believed
 given \(A\).
\end{enumerate}
\end{hbtheorem}

Item~2 is the source's correspondence: for an antecedent that is not
disbelieved, belief in the material conditional and conditional belief
coincide, at every threshold.  In the Tweety prior, \(\kappa(B)=0\) and both
\(\neg B\lor F\) and flight given \(B\) are believed with the same strength
\(1\).  Item~3 is the vacuous case: \(\kappa(P)=1\), so \(\neg P\lor F\)
and \(\neg P\lor\neg F\) are both believed, although only nonflight is
believed given \(P\).  The word ``precisely'' must therefore be read with
the hypothesis \(\kappa(A)=0\).  Conditional beliefs with a disbelieved
antecedent, such as \(\kappa(F\mid P)=3\), still exist, but belief in their
material counterparts carries no information about them.  When
\(0<\kappa(A)\leq z\), neither item~2 nor item~3 applies, and the material
belief at threshold \(z\) holds exactly when
\(\kappa(A)+\kappa(\neg B\mid A)>z\), a mixture of the two attitudes.

\section{Finite firmness, revision and contraction}

\begin{hbdefinition}{def:p15-rat53-soft-update}{Two-cell conditional preservation}
Assume \(\kappa(E),\kappa(\neg E)<\infty\).
Specify new cell ranks \(a,b\in\mathbb N\) with \(\min(a,b)=0\).
Preserve conditional ranks within \(E\) and \(\neg E\), and set
\[
 \kappa^{a,b}_E(A)=
 \min\{\kappa(A\mid E)+a,\ \kappa(A\mid\neg E)+b\}.
\]
Accepting \(E\) with finite firmness \(n>0\) is the case \(a=0,b=n\).
Setting \(a=b=0\) suspends belief between the two cells.  If \(E\)
was previously believed, this gives the source's belief-contraction case;
it should not be called contraction by \(E\) when \(E\) was already
disbelieved, because it would also remove belief in its negation.
For a threshold \(z\), new belief in \(E\) requires \(n>z\).
The parameter is the target firmness of this update, not an unspecified
additive increase of the old two-sided rank.
\end{hbdefinition}

\begin{hbtheorem}{thm:p15-rat53-update}{proved=annex:p15.rat53.update}{Existence and uniqueness of the specified update}
Under the two finite-cell conditions, \(\kappa^{a,b}_E\) is a normalized
negative ranking function, has cell ranks \(a,b\), and preserves both old
within-cell conditional rank functions.  It is the unique ranking function
with those properties.  On finite \(W\), its world ranks are
\[
 r'(w)=
 \begin{cases}
 r(w)-\kappa(E)+a,&w\in E,\\
 r(w)-\kappa(\neg E)+b,&w\notin E.
 \end{cases}
\]
An input with infinite-rank cell is not covered by these two-sided formulas.
Hard conditioning on a finite-rank \(E\) is instead
\(A\mapsto\kappa(A\mid E)\), which gives infinite rank to \(\neg E\).
\end{hbtheorem}

In the table, the first update takes \(E=B,n=2\).  The two old cell ranks
are zero, so it leaves all bird worlds unchanged and raises every nonbird
world by two.  The zero set becomes the flying nonpenguin bird: the belief
set expands.  The next update takes \(E=P,n=1\), now with
\(\kappa'(P)=1\) and \(\kappa'(\neg P)=0\).  It subtracts one on
penguin worlds and adds one on nonpenguin worlds.  The new zero set is the
nonflying penguin bird: some earlier beliefs are replaced.  In particular
\(\tau(P)\) changes from \(-1\) to \(+1\), not by just one unit.
This verifies the source's specified target-firmness example without conflating
it with other proposed update rules
\parencite[p.~340 and n.~6]{KernIsbernerSkovgaardOlsenSpohn2021Ranking}.

The ranking left by an update can be updated again, with new evidence cells
and explicit strengths; one-step belief-set membership alone does not retain
all that information.  On a finite propositional space with all world ranks
finite, acceptance with any \(n>0\) selects the old minimum-rank
\(E\)-worlds.  This is the faithful-preorder selection mechanism of
Definition~\ref{def:p08-agm-faithful-preorder}, with its prior and consistency
hypotheses retained, not a proof that every iterated ranking operation follows
from the eight one-step AGM postulates.

The chapter's cardinal-scale argument points to a separate measurement theorem
of Hild and Spohn: suitably constrained iterated contraction behavior
determines ranks up to a positive multiplicative factor, and its note~7 says
twofold nonvacuous contractions suffice.  It does not print the measurement
axioms or that theorem's proof.  This is a source-cited representation claim,
not a measurement algorithm newly derived here, and not evidence that a
participant's observed responses identify a unique absolute integer scale.
Changing the scale also changes numerical strengths and thresholds; interpersonal
comparison therefore remains an extra issue
\parencite[p.~340 and n.~7]{KernIsbernerSkovgaardOlsenSpohn2021Ranking}.
With the printed integer codomain, a change of units must still give
allowed integer values; permitting arbitrary positive real rescaling is
a real-valued extension, not literally the same codomain.
The learning rule concerns evidence change, not forgetting or every possible
cause of psychological change.  Note~6 links its preservation structure to
Jeffrey conditioning and names Shenoy's alternative; that attribution does
not make all those numerical rules identical.

\section{What the comparisons establish}

The chapter identifies antecedents in Shackle's potential surprise, Rescher's
hypothetical reasoning and Cohen's Baconian probability.  Its claimed new step
is cardinal conditional ranks and a dynamics, not mere ordinal exception
layers.  These are historical attributions made in the chapter; none turns
all referenced books into a new treatment here.

\begin{hbtheorem}{thm:p15-rat53-possibility}{proved=annex:p15.rat53.possibility}{The precise finite static possibility transform}
For finite normalized world ranks set
\[
 \Pi(A)=2^{-\kappa(A)},\quad 2^{-\infty}=0,\qquad
 N(A)=1-2^{-\kappa(\neg A)}.
\]
Then \(\Pi(W)=1\), \(\Pi(\varnothing)=0\), and
\(\Pi(A\lor B)=\max\{\Pi(A),\Pi(B)\}\).
Also \(N(W)=1\), \(N(\varnothing)=0\), and
\(N(A\land B)=\min\{N(A),N(B)\}\).
The inverse on the range \(\{0\}\cup\{2^{-j}:j\in\mathbb N\}\)
recovers the integer ranks.  General real-valued possibility measures are not
all on this grid.  Where \(\kappa(A)<\infty\), the rank conditional
corresponds to \(\Pi(A\land B)/\Pi(A)\), not automatically to every
ordinal possibility-conditioning convention.
\end{hbtheorem}

Thus the source's possibility-equivalence claim has a precise static
representation sense, with range and interpretation qualifications.
The distinct Dubois--Prade treatment in
Chapter~\ref{sec:p15-rationality-4-7}, particularly
Theorem~\ref{thm:p15-r47-ranking}, discusses possibility, necessity and
belief functions; it does not supply the present rank-arithmetic dynamics.
Consonant belief functions can represent the same finite static necessity
values, but a common representation does not by itself select the same
learning operation or establish betting-price semantics
\parencite[pp.~341--342]{KernIsbernerSkovgaardOlsenSpohn2021Ranking}.

The probability analogy uses negative ranks: sum over alternatives becomes
minimum, multiplication becomes addition, and division becomes subtraction.
It is an analogy of operations, not an equality of numbers or interpretations.
For example two disjoint finite events can both have zero rank, whereas two
disjoint events cannot both have probability one.  As an editorial illustration,
for a finite regular world ranking the normalized family
\(P_\varepsilon(w)=\varepsilon^{r(w)}/\sum_v\varepsilon^{r(v)}\)
has rank \(\kappa(A)\) as the limiting exponent
\(\lim_{\varepsilon\downarrow0}\log P_\varepsilon(A)/\log\varepsilon\)
for nonempty \(A\).  The annex proves this restricted bridge; neither a
single ordinary probability nor a universal fixed acceptance threshold is
thereby recovered.  The source's disputed bridge between probability and
outright belief is compared with the probability norms of
Chapter~\ref{sec:p15-rationality-4-1}, not settled by replacing its credences
with ranks.

The same operation analogy underlies the chapter's remark that the theory of
Bayes nets could equally well be stated in ranking terms, for which it cites
Spohn's 1988 paper and Goldszmidt and Pearl
\parencite[p.~340]{KernIsbernerSkovgaardOlsenSpohn2021Ranking}.
Chapter~\ref{sec:p16-rationality-4-2} presents the probabilistic theory: the
parental Markov condition of Definition~\ref{def:p16-rat42-bayes-net}, the
product representation of Theorem~\ref{thm:p16-rat42-product}, and the
conditional-independence operations of Lemma~\ref{thm:p16-rat42-semigraphoid}.
The step that transfers directly is the chain rule.

\begin{hbtheorem}{thm:p15-w02b-p15-039-chain}{proved=annex:p15-w02b-p15-039-chain}{The additive chain rule for ranks}
Let \(C_1,\ldots,C_n\) be propositions, \(n\geq2\), with
\(\kappa(C_1\land\cdots\land C_{n-1})<\infty\).  Then
\[
 \kappa(C_1\land\cdots\land C_n)=\kappa(C_1)+\sum_{i=2}^{n}
 \kappa(C_i\mid C_1\land\cdots\land C_{i-1}),
\]
every conditional rank on the right being defined.  For a finite partition
\(\{D_1,\ldots,D_m\}\) of \(W\) and any \(A\),
\(\kappa(A)=\min_j\bigl(\kappa(D_j)+\kappa(A\mid D_j)\bigr)\), the minimum
being taken over the cells with \(\kappa(D_j)<\infty\).
\end{hbtheorem}

Thus a ranking network replaces the product of local tables by a sum of local
conditional ranks and the summation over unobserved values by a minimum,
exactly as the source's translation of operations prescribes.  The
probabilistic product theorem in Chapter~\ref{sec:p16-rationality-4-2} also
needs the parental Markov condition, a statement about conditional
independence.  The ranking counterpart of that condition, and the ranking
versions of the semigraphoid operations, are attributed by the source to
the cited works; the chapter states neither, and they are not derived here.
The present chain rule is therefore the arithmetic half of the restatement
claim, not a proof that every graphical independence theorem holds for ranks.

\section{AI methods and psychological standing}

The chapter explains why rankings supply a common computational state for
nonmonotonic consequence and iterated revision: a conditional may be evaluated
at minimum-ranked antecedent worlds, and an update leaves another ranked
state.  This connects to the stronger ranked-model representation of
Theorem~\ref{thm:p08-rational-representation}, not to every smooth preferential
model.  Its account of System Z, the Darwiche--Pearl tradition and later
multiple-change approaches is a survey of applications, not one interchangeable
calculus.  Its conditional-preservation discussion distinguishes c-revisions
from c-representations obtained from a uniform initial rank.  Preservation
of within-cell differences is exhibited explicitly above; no unprinted full
postulate system or general simultaneous-conditional algorithm is attributed
to this short chapter \parencite[pp.~341--342]{KernIsbernerSkovgaardOlsenSpohn2021Ranking}.

For a finite world representation, \(A\mathrel{|\!\sim}B\) means that
the minimum-ranked \(A\)-worlds satisfy \(B\).  If
\(\kappa(A)<\infty\), this is
\(\kappa(A\land\neg B)>\kappa(A)\).
The computation requires a representation of those worlds or an appropriate
symbolic algorithm; a vocabulary with many atoms can have exponentially
many worlds.  Simpler min/add arithmetic alone does not prove a general
complexity advantage over all probability algorithms.

For an arbitrary negative ranking function and \(\kappa(A)<\infty\), write
\(A\mathrel{|\!\sim}_\kappa B\) when \(\kappa(A\land\neg B)>\kappa(A)\);
on a finite world representation this is the minimum-world reading just given.
The chapter's psychological section reports that Da Silva Neves, Bonnefon and
Raufaste tested whether participants' possibility judgments satisfy the
postulates of System P, augmented by rational monotonicity: from
\(A\mathrel{|\!\sim}C\) and not \(A\mathrel{|\!\sim}\neg B\), infer
\(A\land B\mathrel{|\!\sim}C\), in cases where responses violated plain
monotonicity, from \(A\mathrel{|\!\sim}C\) to \(A\land B\mathrel{|\!\sim}C\)
\parencite[p.~342 and p.~344, n.~9]{KernIsbernerSkovgaardOlsenSpohn2021Ranking}.
Ranking consequence satisfies the added postulate directly.

\begin{hbtheorem}{thm:p15-w02b-p15-039-rational-monotony}{proved=annex:p15-w02b-p15-039-rational-monotony}{Rational monotonicity of ranking consequence and the Tweety failure of monotonicity}
Let \(\kappa(A)<\infty\).  If \(A\mathrel{|\!\sim}_\kappa C\) and not
\(A\mathrel{|\!\sim}_\kappa\neg B\), then \(\kappa(A\land B)=\kappa(A)\)
and \(A\land B\mathrel{|\!\sim}_\kappa C\).
In the Tweety prior, \(B\mathrel{|\!\sim}_\kappa F\) and
\(B\land P\mathrel{|\!\sim}_\kappa\neg F\), so monotonicity fails; the
rational-monotonicity premise is not met there, because
\(B\mathrel{|\!\sim}_\kappa\neg P\).
\end{hbtheorem}

This gives the comparison with Chapter~\ref{sec:p15-rationality-4-4}.  There
System P is the reference family for default inference, read through coherent
conditional probability, and unrestricted strengthening of the antecedent is
excluded.  Rational monotonicity is a further postulate, not one of System P's
rules; it is the condition that the ranked-model representation of
Theorem~\ref{thm:p08-rational-representation} adds to System P.  A test of
System P together with rational monotonicity is therefore a test of the ranked
structure that possibility and ranking functions share, not of probability
logic's interval bounds.  The chapter stresses the limit of that evidence:
these were possibility judgments, and no corresponding direct test of ranking
theory using judgments of disbelief or implausibility had been made.  The
chapter gives no sample, task-by-task result or uncertainty estimate for the
Da Silva Neves study; none is supplied here, and nothing about the size or
generality of its findings follows from this summary.  Chapter~\ref{sec:p15-rationality-4-4}
raises a further objection that applies equally here: agreement of responses
with a rule system does not by itself establish that rule system's normative
authority, nor how participants interpreted the task.

The psychological section identifies possible uses, indirect tests and
obstacles.  Its possibility-judgment studies investigate defeasible
inference patterns, not direct measurements of negative-rank values.
Its language-comprehension application models the processing cost of
implausible information, while its suppression discussion explicitly says
that ranking machinery alone does not produce the Lisa/library suppression
effect: assumptions about activated background knowledge are required.
The chapter also discusses "normally" in MP/MT tasks, reason-relation and
relevance evaluations, and individual variation.  These are attributed
research programmes, not a claim that one ranking model has been directly
validated across all their populations
\parencite[pp.~342--344]{KernIsbernerSkovgaardOlsenSpohn2021Ranking}.

The chapter's first psychological proposal concerns the belief revision
research reviewed in Chapter~\ref{sec:p19-rationality-5-4}, which it describes
as inspired by AGM and by probabilistic updating.  It suggests that such work
could be extended by ranking theory, because ranking theory was meant to
represent full belief, unlike probabilistic updating, while improving on AGM
in allowing iterated revision
\parencite[p.~342]{KernIsbernerSkovgaardOlsenSpohn2021Ranking}.
Two results of that chapter locate exactly what the extension would and would
not add.  Theorem~\ref{thm:p19-r54-two-revisions} gives two total orders on
the valuations of \(p,q\), both with the belief set
\(\operatorname{Cn}(\{p\to q,p\})\), whose revisions by \(\neg q\) keep
opposite premises.  Ranks realizing them are
\[
\begin{aligned}
 \kappa_1&:\ (1,1)\mapsto0,\ (1,0)\mapsto1,\ (0,0)\mapsto2,\ (0,1)\mapsto3;\\
 \kappa_2&:\ (1,1)\mapsto0,\ (0,0)\mapsto1,\ (1,0)\mapsto2,\ (0,1)\mapsto3.
\end{aligned}
\]
In both, \(\kappa(\neg q)=1\), so the contrary fact is disbelieved with
finite firmness and conditioning on it is defined by
Definition~\ref{def:p15-rat53-conditionals}.  The posterior zero-rank world is
\((1,0)\) under \(\kappa_1\), keeping \(p\), and \((0,0)\) under
\(\kappa_2\), keeping \(p\to q\); either posterior is again a ranking state
available for a further change.  By contrast,
Theorem~\ref{thm:p19-r54-probability-boundary} records that ordinary
conditioning needs positive prior probability, so a probability assigning the
believed \(q\) probability \(1\) cannot be conditioned on \(\neg q\).  Ranking
theory thus supplies what its motivation promised: an outright belief that is
revisable by conditioning and leaves an iterable state.  It does not decide
which premise a participant keeps.  That choice is fixed by the prior ranks,
and ranks have not been elicited in the experiments reviewed there; the
three obstacles with which the chapter closes, listed below, concern
precisely that gap.  The
psychology chapter's distinction between withdrawing a derived conclusion
and revising a premise also remains: the ranking update acts on a whole
state over valuations, while the laboratory response is a choice between
two sentences.

The source's suppression illustration gives Lisa an essay to write and the
conditional that if she has an essay she will study late in the library.
It then adds the conditional that if the library is open she will study late.
As material implications, these premises still entail late study by modus
ponens; the added premise cannot invalidate that derivation.  A defeasible
model must instead represent the activated possible defeater, such as the
library's being closed.  That background constraint is not logically implied
by the added open-library conditional alone.  The chapter's point is precisely
that the bare ranking inference mechanism does not generate the suppression
pattern without further knowledge-base assumptions; endorsement behavior,
conditional interpretation and deductive validity must not be conflated.

One directly inspectable study is Skovgaard-Olsen, Singmann and Klauer's
2016 relevance experiment.  Its corrected accepted manuscript reports
577 completed online sessions and 348 included native-English participants
from the USA, UK and Australia, after comprehension, timing and seriousness
exclusions.  It manipulated positive, negative or irrelevant antecedent--
consequent relationships across realistic scenarios, and asked for probability
or acceptability of indicative or concessive conditionals.  In the indicative
groups the estimated slope relating conditional probability to the response
was \(0.78\) (95\% CI \([0.71,0.85]\)) with positive relevance,
\(0.61\) (\([0.52,0.71]\)) with negative relevance and
\(0.43\) (\([0.35,0.52]\)) with irrelevance.  This is evidence about
that task's relevance modulation, not identification of a negative ranking
function or a universal semantics for every conditional.  Background-knowledge
manipulations, participant exclusions and the competing conditional-probability
account belong to its interpretation
\parencite[accepted-manuscript pp.~16--23]{SkovgaardOlsenSingmannKlauer2016Relevance}.

Juslin and colleagues' 2011 log-format work concerns probability judgment,
not a direct rank elicitation.  Its first experiment used 30 undergraduates
in probability, odds and log-odds versions of base-rate problems.  Format
varied within participants and metric versus computational instruction
between participants: this was not a pre/post instruction design.
Each participant solved 54 problems; performance was measured by mean
absolute error against the corresponding Bayesian answer.  The study reports
better error performance in the log-odds format, but explicitly calls such
bias reductions nominal: they need not imply conceptual understanding or
generalization to novel tasks.  Such format effects provide
a reason to investigate additive representations, not a deduction that
integer disbelief ranks are the mechanism used.  The chapter presents this
as motivation, alongside the lack of direct rank-scale tests
\parencite[pp.~253--257]{JuslinEtAl2011Logarithms}.

Three obstacles conclude the chapter.  First, negative ranks require thinking
about disbelief in negations; switching to two-sided ranks introduces distinct
positive and negative scale regions.  Second, the contraction-based scale has
not been experimentally implemented like betting quotients.  Third, an
unbounded integer scale has no canonical decimal bins for participants.
The source's 2014 model and 2017/2019
reason-relation work are linked proposals and findings whose specific
protocols cannot be inferred from their bibliographic titles.  The normative
ranking postulates, a fitted conditional model and direct measurement of
disbelief are separate evidential targets.

\begin{hbmethod}{def:p15-rat53-task-interface}{Choose a rank representation and a learning operation}
Specify the possibility space, proposition algebra, rank range and whether
finite or arbitrary belief closure is claimed.  Declare the threshold for
outright belief.  For conditioning give the prior rank of the evidence; an
infinite value blocks the subtraction rule.  For uncertain learning give both
new cell ranks and which old within-cell conditionals are retained.  A phrase
such as "new evidence" does not determine its firmness.

To compare with probabilities distinguish static transforms, limiting
families, update rules and elicitation interpretations.  To reason with a
conditional distinguish material implication, conditional belief and an
increase relative to the opposite antecedent.  To make a psychological claim
state the actual task, sample, measured response and model comparison; Tweety's
consistent arithmetic is not an experiment.
\end{hbmethod}
